% Grid-current clipping, blocking and recovery (triode-overload): the set-up a person would wire, semantically the deck's.
%   VIN in 0 PWL(<burst>), CIN in x <cin>, RSTOP x gj <rstop>, RGL gj 0 1meg, VGM gj g 0   the burst through the coupling capacitor and the stopper; the grid leak at the grid
%   the rest of the stage as the RC chart's middle cell, Rk with its bypass capacitor
%   .tran over the burst and 200 ms after it
\documentclass[tikz, border=10pt]{standalone}
\usepackage{bench}
\begin{document}
\begin{tikzpicture}
\benchtitle{0,8.6}{Grid-current clipping, blocking and recovery}
\def\gnd{1.0}
% the triode
\node[triode, scale=1.35] (T) at (5.6,3.8) {};
\node[anchor=west, font=\small, align=left] at ($(T.anode)!0.62!(T.cathode)+(0.75,0)$) {DUT};
% copy : the anode resistor up to the rail, the anode ammeter under it
\coordinate (am) at ($(T.anode)+(0,1.1)$);
\ammeter{T.anode}{am}{$I_a$}
\draw (am) to[R, l_={$R_p$}] ++(0,1.6) coordinate (top);
\coordinate (vtop) at ($(top)+(6.6,0)$);
\draw (top) -- (vtop);
\draw[set] (vtop) to[V, l={$E_{bb}$}] (vtop |- 0,\gnd) node[ground, black]{};
\nodename{anchor=south}{$(top)+(1.0,0.05)$}{bb}
\nodename{anchor=east}{$(am)+(-0.15,0.15)$}{a}
% the cathode resistor, and its bypass
\draw (T.cathode) to[R, l={$R_k$}] (T.cathode |- 0,\gnd) node[ground]{};
\draw (T.cathode) -- ++(1.2,0) coordinate (ck) to[C, l={$C_k$}] (ck |- 0,\gnd) node[ground]{};
\nodename{anchor=east}{$(T.cathode)+(-0.1,-0.25)$}{k}
% the grid: its resistor to ground and the source
\coordinate (g) at (T.control);
\draw (g) -- ++(-0.6,0) coordinate (gj);
\draw (gj) to[R, l_={$R_{gl}$}] (gj |- 0,\gnd) node[ground]{};
\draw (gj) -- ++(-0.3,0) coordinate (c0);
\draw (c0) to[R, l_={$R_{stop}$}] ++(-1.6,0) coordinate (c1);
\draw (c1) to[C, l_={$C_{in}$}] ++(-1.6,0) coordinate (c2);
\coordinate (vin) at (c2);
\draw[set] (vin) to[V, l_={$v_{in}$}] (vin |- 0,\gnd) node[ground, black]{};
\nodename{anchor=south}{$(g)+(-0.2,0.05)$}{g}
\nodename{anchor=south east}{$(vin)+(-0.05,0.05)$}{in}
% the output: through Cc into the next grid's resistor
\coordinate (oj) at ($(am)+(0,0.0)$);
\draw (am) -- ++(1.6,0) to[C, l={$C_c$}] ++(1.4,0) coordinate (o);
\draw (o) to[R, l_={$R_s$}] (o |- 0,\gnd) node[ground]{};
\nodename{anchor=south}{$(o)+(0.3,0.05)$}{o}
% the voltmeter on o
\coordinate (vo) at ($(o)+(1.2,0)$);
\draw (o) -- (vo);
\voltmeter{vo |- 0,\gnd}{vo}{$v_{out}$}
\draw (vo |- 0,\gnd) node[ground]{};
% the definitions
\definitions{0,-0.3}{%
  a 1 kHz burst of 20 periods at 1, 3, 10 and 30 times the drive at which grid current starts, then 1/10 of that drive\\[2pt]
  $C_{in} = 22$ nF;\quad $R_{gl} = 1$ M\ohm;\quad $R_{stop} = 0$, 1 k, 68 k\ohm;\quad $C_k = 25$ \textmu F\\[2pt]
  read the grid's mean voltage during and after the burst (the bias the grid current charges $C_{in}$ to), the gain after it, and the time for the gain to come back within 10\,\%\\[2pt]
  a model without grid current neither clips at the grid nor blocks: `absent', and the guitar-amplifier behaviour it stands for with it%ENDDEFS
}
\end{tikzpicture}
\end{document}
